% III. Проектиране на системата.
% Слоеве, модел на данните, формати. Диаграмите се добавят, когато слоевете са
% фиксирани. Празна фигура е по-лоша от липсваща.
\section{Проектиране на системата}
\label{sec:design}

\subsection{Слоеве и отговорности}

Системата е разделена на четири слоя с еднопосочна зависимост отгоре надолу.
Слоят на \term{модела} съдържа описанието на бизнес процеса и на проектите
кандидати и не зависи от нищо друго. Слоят на \term{оценяването} превръща
описанието в число за всеки проект чрез избрания метод за многокритериална
оценка. Слоят на \term{подбора} формулира и решава задачата за избор на
портфейл. Слоят на \term{представянето} приема сценарий, извиква останалите три
и връща резултата заедно с обяснението към него.

Границите между слоевете са поставени така, че смяната на метода за оценяване да
не засяга подбора, а смяната на алгоритъма за подбор да не засяга оценяването.
Това не е абстракция за всеки случай, а пряко следствие от
Раздел~\ref{sec:alternatives}: и двата слоя имат по повече от една реализация още
в първата версия. Потокът от данни между слоевете е показан на
Фиг.~\ref{fig:layers}.

\begin{figure}[H]
\centering
\begin{tikzpicture}[node distance=0.75cm and 1.1cm]
  \node[block] (scen) {Сценарий\\\code{scenarios/*.json}};
  \node[block, below=of scen] (model) {Слой на модела\\\code{model.cpp}, \code{cpm.cpp}\\граф, инварианти,\\критичен път, PERT};
  \node[block, below=of model] (score) {Слой на оценяването\\\code{mcdm.cpp}\\аналитична йерархия, TOPSIS};
  \node[block, below=of score] (select) {Слой на подбора\\\code{portfolio.cpp}};
  \node[accent, left=of select] (greedy) {Лакома евристика\\долна граница};
  \node[accent, right=of select] (exact) {Разклоняване и\\ограничаване\\линейна релаксация};
  \node[danger, below=of select] (check) {Проверка:\\точно $\ge$ евристика};
  \node[block, below=of check] (post) {Риск и чувствителност\\\code{risk.cpp}, \code{sensitivity.cpp}};
  \node[block, below=of post] (out) {Представяне\\\code{report.cpp}, \code{main.cpp}};

  \draw[line] (scen) -- (model);
  \draw[line] (model) -- (score);
  \draw[line] (score) -- (select);
  \draw[line] (select) -- (greedy);
  \draw[line] (select) -- (exact);
  \draw[line] (greedy) |- (check);
  \draw[line] (exact) |- (check);
  \draw[line] (check) -- (post);
  \draw[line] (post) -- (out);
\end{tikzpicture}
\caption{Слоеве и поток от данни. Стрелките сочат посоката на зависимостта.}
\label{fig:layers}
\end{figure}

Проверката преди последните два блока не е украса на диаграмата. Тя е
единственото място, където грешка в построяването на модела може да бъде уловена
автоматично: ако точният алгоритъм върне стойност под тази на допустимо решение,
намерено от евристиката, изпълнението спира със съобщение за грешка в модела,
вместо да отчете число.

\subsection{Модел на данните}

Бизнес процесът се представя като насочен ацикличен граф. Върховете са дейности
с продължителност, стойност и потребен ресурс. Ребрата са отношения на
предхождане. Проектът кандидат е свързано подмножество от дейности заедно с
оценките си по критериите. Портфейлът е подмножество от проекти.

Ацикличността е инвариант, който системата проверява при зареждане, а не
допускане. Цикъл в графа на предхождането прави най-ранният срок на завършване
недефиниран, а задачата за подбор неразрешима; откриването му при зареждането
дава на потребителя посочимо място в неговия вход, вместо съобщение за грешка от
алгоритъма. Проверката се прави на двете нива поотделно: между дейностите на
един проект и между самите проекти.

\begin{table}[H]
\centering
\caption{Обекти на модела и техните задължителни полета}
\label{tab:model}
\begin{tabular}{@{}L{3.0cm}L{4.4cm}L{4.6cm}@{}}
\toprule
\textbf{Обект} & \textbf{Задължителни полета} & \textbf{Инвариант} \\
\midrule
Дейност & идентификатор, продължителност, стойност & неотрицателни стойности, оптимистична $\le$ най-вероятна $\le$ песимистична \\
Зависимост & източник, цел & графът остава ацикличен, и двата края съществуват \\
Проект & идентификатор, дейности, оценки по критериите & непразно множество дейности, оценка по всеки критерий \\
Критерий & идентификатор, посока & сумата на теглата се нормира до единица \\
Ресурс & идентификатор, капацитет & неотрицателен капацитет \\
Ограничение & бюджет, капацитет по ресурс, предхождане, взаимно изключване & предхождането и изключването не се отнасят за една и съща двойка \\
\bottomrule
\end{tabular}
\end{table}

Съставът на модела е обобщен в Табл.~\ref{tab:model}. Посоката на критерия указва
дали по-голямата стойност е по-добра, което е необходимо и за нормализацията в
TOPSIS, и за обръщането на колоната при синтеза по аналитичната йерархия.

Три величини не се въвеждат, а се извеждат при зареждането, за да няма два
източника на едно и също число. Стойността на проекта е сумата от очакваните
стойности на дейностите му. Продължителността на проекта е дължината на неговия
критичен път. Потребният ресурс на проекта е върховото едновременно потребление
при разписание по най-ранни срокове, изчислено чрез обхождане на моментите на
започване и завършване. Последното е съзнателно консервативно: то допуска, че
всеки проект работи с върховото си потребление през целия период, в който е
активен, и следователно никога не подценява натоварването.

\subsection{Формат на сценария и на резултата}

Входът е JSON документ с фиксирана схема. Сценарият съдържа критериите, матрицата
на попарните сравнения, ресурсите с техните капацитети, бюджета и проектите с
дейностите и оценките им. Пълната схема е дадена в Приложение~А.

Изходът е текстов доклад с таблици, а не втори JSON документ. Причината е, че
изходът не се подава на друга програма: той се чете от човек и се прилага към
решението. За всеки отхвърлен проект се отпечатва причината за отхвърлянето, като
причините се търсят в определен ред: непокрито предхождане, взаимно изключване с
финансиран проект, недостиг на бюджет, недостиг на ресурс. Ако нито едно от тези
не важи, проектът е бил допустим и е отпаднал по стойност, и това се казва
изрично, вместо да се посочи ограничение, което не е било определящо.

\subsection{Анализ на риска}

Тройните оценки в сценария носят несигурност, която очакваната стойност изхвърля.
Анализът на риска я връща обратно, като извлича извадка от разпределението на
всяка продължителност и всяка стойност и преизчислява разписанието и сумарната
стойност за всяка извадка. Резултатът е разпределение на общата стойност на
портфейла и на срока на завършване на програмата, от което се отчитат средна
стойност, медиана и централен интервал.

Проектното решение тук е разпределението да бъде триъгълно върху трите точки.
То се определя изцяло от числата, които сценарият вече съдържа, няма
допълнителен параметър за форма, който да бъде предположен, и обратната му
функция е в затворен вид, което при стотици хиляди извадки има значение.
Разликата спрямо бета разпределението, което стои зад очакваната стойност на
PERT, е отчетена в Раздел~\ref{sec:implementation}.

\subsection{Анализ на чувствителността}

Анализът на чувствителността се проектира като повторно изпълнение на същия
модел върху систематично изменен вход, а не като отделен алгоритъм. Изменят се
теглата на критериите и бюджетът, всяко в зададен интервал и стъпка. Теглата се
изменят по едно наведнъж и след това се нормират наново, така че всяко изпълнение
да се отнася до едно-единствено изменено допускане; едновременното изменение би
дало резултат, който не може да бъде приписан на определена причина.

Резултатът е множество от портфейли, от които се извежда кои позиции остават
избрани при всички изменения и кои се появяват и изчезват. Първите са
устойчивата част от препоръката и това е основният изход, предназначен за
управленския читател. Освен обобщението по позиции се отпечатва и решетка с по
един стълб на изпълнение, защото делът сам по себе си казва колко често една
позиция е била финансирана, но не и при какви допускания.
